\documentclass[11pt]{article}

% to avoid loading the natbib package, add option nonatbib:
%\usepackage[nonatbib]{neurips_2025}
\usepackage{iclr2025_conference,times}

\iclrfinalcopy %comment for anonymous submission

\usepackage[australian]{babel}
\usepackage[T1]{fontenc} %For international vectorised fonts
\usepackage{amssymb}
\usepackage{amsmath}
\usepackage{amsthm}
\usepackage{mathrsfs} % For Cursive Fonts, dont define with mathdesign
\usepackage{textcomp} %Text symbols

\usepackage[breaklinks,plainpages=false,pdfpagelabels]{hyperref}
\hypersetup{colorlinks,citecolor=blue,filecolor=blue,linkcolor=blue,urlcolor=blue}
%\hypersetup{colorlinks,citecolor=black,filecolor=black,linkcolor=black,urlcolor=black}

%\usepackage{geometry}
\usepackage{graphicx}
\usepackage[caption=false, font=footnotesize]{subfig}
\usepackage{nicefrac}
\usepackage{fixmath}
\usepackage{bm}
\usepackage{array}
\usepackage{mdwmath} %Math alignment
\usepackage{mdwtab}
\usepackage[printonlyused,withpage]{acronym} %acronyms package
\usepackage{microtype} %Cleans up text width
\usepackage{marginnote}

%-----------------Theorems-----------------
\newtheorem{theorem}{Proposition}
\newtheorem{corollary}{Corollary}
\newtheorem{definition}{Definition}
\renewcommand{\qedsymbol}{$\blacktriangleleft$}

%-----------------Math-----------------
\newcommand{\mat}[1]{\mathbold{#1}} % matrix style 
%Better Modulo notation imod 
\makeatletter
\def\imod#1{\allowbreak\mkern10mu({\operator@font mod}\,#1)}
\makeatother
\newcommand{\wholeNumbers}{\mathbb{N}}
\newcommand{\rationalNumbers}{\mathbb{Q}}

%-----------------Defines-----------------
\makeatletter
\renewcommand*{\fps@figure}{htb}
\renewcommand*{\fps@table}{htb}
\makeatother

% Slightly bigger margins than the latex defaults
%\geometry{verbose,tmargin=1in,bmargin=1in,lmargin=0.5in,rmargin=0.5in}

%Title of paper
%\title{Neural circuits with self-engineered cellular machines}
\title{Neural networks as self-engineered cellular machines}
%\title{Self-engineered neural networks via learnable cellular machines} 
%\title{Self-engineered networks with learnable neural circuits} 
%\title{Neural self-engineering of networks with learnable cellular circuits} 

\author{Shekhar S. Chandra\\
	%\thanks{shekhar.chandra@uq.edu.au}
	%\cortext[cor1]{Principal corresponding author}
	School of Electrical Engineering and Computer Science,\\
	The University of Queensland, Australia.\\
	\texttt{shekhar.chandra@uq.edu.au} }

%=================Start Doc===============
\begin{document}
\maketitle

\begin{abstract}
	
\end{abstract}
	
\section{Introduction}\label{sec::Introduction}
%In this work, we unite the modern deep learning framework to solve tasks through neural networks with cellular automata that use simple rules to assemble structures. The resulting framework allows neural networks to additionally learn their architecture and cellular automata to learn the rules for their assembly.
%Artificial neural networks have become very successful recently through the advancements in deep learning.

In the mid-twentieth century, John von Neumann envisioned a simple neural system to mimic the human brain comprising cells of neurons, each of which maintaining a state or purpose and used it to construct a theory of self-replicating machines~\citep[published post-humorously]{von_neumann_theory_1966}. This work resulted in one of the founding works of cellular automata, which is a simulation within a lattice of cells that evolve via simple rules based on the current status of each cell. He envisioned that the cells in these cellular systems would approximate neurons that behave much like the McCulloch and Pits formulation of an artificial neuron devised around the same time~\citep{mcculloch_logical_1943}. His key insight was that these cells would make decisions like neurons in the brain but would also maintain states depending on their purpose~\citep{neumann_computer_1958}.
%proposed machines as . These cells have different specialised roles depending on the task, where they were originally introduced during machine replication, a process that resulted in the first independent construction of cellular automata~\citep[published post-humorously]{von_neumann_theory_1966}. He envisioned that machines such as these consist of cells that approximate neurons that behave much like the McCulloch and Pits formulation of an artificial neuron devised around the same time~\citep{mcculloch_logical_1943}. His key insight was that these cells would make decisions like neurons in the brain but would also maintain states depending on their purpose~\citep{neumann_computer_1958}. 
However, the majority of advancements since von Neumann's original proposal has been largely focused on simplifying the rules and states required of the cellular automata in order to produce Turing complete systems; see a brief review of this area provided in Appendix~\ref{sec::review-replication}.
%this original intent of von Neumann in using neurons for self-replication and cellular automata has been largely unexplored since it was originally proposed in favour of simpler designs; see a brief review of this area provided in~\ref{sec::review-replication}.

In recent years, the area of deep learning has solved many problems associated with artificial neural networks and has been used to solve a large number of important problems~\citep{lecun_deep_2015}.
% neurons improvements
Neurons have evolved from~\citet{mcculloch_logical_1943} style designs to the modern definition of how a single neuron activates as
\begin{equation}
    \bar{x} = \sigma (W\cdot x+b),
\end{equation}
where $x$ is the input, $W$ represents the weights of the various (usually fixed) connections to this neuron, $b$ is the bias of the neuron and $\sigma$ is an activation function applied when firing.
% mixing mechanisms
% architectures
% back prop
Various attempts have been made to add memory and states to deep learning, but these have had largely unstable training, although they were able to learn to solve simple problems~\citep{graves_neural_2014, graves_hybrid_2016}. A few attempts have been made to produce learnable cellular automata including CoDI~\citep{gers_codi-1bit_1998} that relies on spiking networks and genetic encoding to simulate its network.

In this work, we create a framework for deep learning as cellular automata by imbuing an artificial neuron with a learnable (continuous) state as John von Neumann originally intended. This allows us to replicate deep learning entirely on cells, resulting in artificial neural networks running as a cellular automata, which we will refer to as~\acp{VNN}.

\appendix
\section{Review of Self-replication}\label{sec::review-replication}
John von Neumann used 29 different states in his original formulation to imbue his cellular neurons with special functions during the self-replication process of his machine, where the machine encodes instructions of a universal Turing machine that eventually computes the replication~\citep{von_neumann_theory_1966}. Significant progress has been made in implementing~\citep{mange_macroscopic_2004,buckley_signal_2008} and extending self-replicating machines of von Neumann in the decades following his work. \citet{codd_cellular_1968}, with subsequent recent amendments from~\citet{hutton_codds_2010}, reduced the requirement of self-replication requiring 29 states to a total of only 8 states. Others have further simplified self-replication to fewer states, reduced the dimensionality, and removed universal requirements~\citep{banks_information_1971,wolfram_statistical_1983,langton_self-reproduction_1984}. A more detailed review of this area can be found in~\citet{sipper_fifty_1998}.

More recent work has focused on simpler cellular automata with as few rules as possible, while maintaining complex behaviour. The most famous of these is John Conway's Game of Life~\citep{gardner_mathematical_1970}, where the cells have only 2 states and rules that change the state of cells based on its neighbours. Surprisingly, even this simple 2-state cellular system is Turing complete~\citep{rendell_fully_2013}. More complex rules have been developed to simulate artificial life~\citep{hong_kong_lenia_2019} and are the subject of ongoing research~\citep{chan_lenia_2020}.

\bibliographystyle{elsarticle-harv} %elsarticle-num
\bibliography{references}

\section{Abbreviations}
\begin{acronym}[GCD]
\acro{1D}{one dimensional}
\acro{2D}{two dimensional}
\acro{3D}{three dimensional}
\acro{VNN}{Von Neumann Neural Network}
\acro{GCD}{greatest common divisor}
\acro{ACE}{automatic computing engine}
\acro{FSM}{Finite State Machine}
\acro{UTM}{Universal Turing Machine}
\acro{UC}{Universal Computation}
\acro{UR}{Universal Replication}
\acro{CPU}{central processing unit}
\acro{GPU}{graphical processing unit}
\acro{ALU}{arithmetic logic unit}
\acro{GoL}{Game of Life}
\acro{QM}{quantum mechanics}
\end{acronym}

\end{document}
